\subsection{内部坐标束屏幕的后验超立方体、外边界精确化与坐标价值原子}\label{subsec:conclusion-screen-hypercube-boundary-torsor-shapley}

沿用定理 \ref{thm:conclusion-screen-entropy-audit-identity}、
定理 \ref{thm:conclusion-coordinate-bundle-kernel-slab-decomposition}、
定理 \ref{thm:conclusion-coordinate-bundle-logmodularity-visible-rank}、
定理 \ref{thm:conclusion-relative-linear-holography-exact-supplement-dimension}、
定理 \ref{thm:conclusion-screen-kolmogorov-deficit-exact-splitting}、
推论 \ref{cor:conclusion-screen-max-conditional-complexity-equals-rank} 与
定理 \ref{thm:spg-internal-coordinate-bundle-screen-cost-closed-form}
的记号。固定
\[
J\subseteq [n],
\qquad
s:=|J|,
\qquad
L_J:=2^{m(n-s)},
\]
并记内部坐标束屏幕为 \(S_J^{\mathrm{int}}\)。令
\[
T_J:=f_{S_J^{\mathrm{int}}}:C_n(\mathcal Q_m;\FF_2)\to \FF_2^{\,S_J^{\mathrm{int}}},
\qquad
Y_J:=T_J(X),
\]
其中 \(X\) 在 \(C_n(\mathcal Q_m;\FF_2)\) 上服从均匀分布。对每个补坐标
\[
a_{J^c}\in\{0,\dots,2^m-1\}^{J^c},
\]
记对应板层为 \(\Sigma_{a_{J^c}}\)，其顶维示性链记为 \([\Sigma_{a_{J^c}}]\)。前述链条已经给出：内部坐标束屏幕的核由板层示性链显式张成，其亏秩、最小审计位数、线性补全维数与最大条件 Kolmogorov 容量都恰等于 \(L_J\)。由此还可进一步抽出一条更细的终端链：后验纤维在每个可见综合征上精确分裂为 \(L_J\) 维独立超立方体；外边界精确化不再是一般拟阵优化，而是按板层完全分离的 hitting 问题；所有最优线性补全构成一个显式仿射 torsor；而隐藏量在坐标方向上的边际价值最终退化为布尔格上的纯一阶 Möbius 原子。

\begin{theorem}[内部坐标束屏幕后验的超立方体因子化]\label{thm:conclusion-coordinate-bundle-posterior-hypercube-factorization}
\leanverified{paper\_conclusion\_coordinate\_bundle\_posterior\_hypercube\_factorization}
对每个 \(y\in \operatorname{im}T_J\)，存在可计算截面
\[
s_J:\operatorname{im}T_J\to C_n(\mathcal Q_m;\FF_2)
\]
使得条件随机变量 \(X\mid (Y_J=y)\) 具有精确表示
\[
X
=
s_J(y)+
\sum_{a_{J^c}\in\{0,\dots,2^m-1\}^{J^c}}
\varepsilon_{a_{J^c}}[\Sigma_{a_{J^c}}],
\]
其中
\[
\bigl(\varepsilon_{a_{J^c}}\bigr)_{a_{J^c}\in\{0,\dots,2^m-1\}^{J^c}}
\]
是 \(L_J\) 个独立公平 Bernoulli 比特。特别地，
\[
X\mid (Y_J=y)\cong \FF_2^{L_J}
\]
作为均匀仿射超立方体。
\end{theorem}
\begin{proof}
由定理 \ref{thm:conclusion-coordinate-bundle-kernel-slab-decomposition}，
\[
\ker T_J
\cong
\bigoplus_{a_{J^c}\in\{0,\dots,2^m-1\}^{J^c}}\FF_2\,[\Sigma_{a_{J^c}}],
\]
故板层示性链在 \(\FF_2\) 上给出 \(\ker T_J\) 的一组基，而
\[
\dim_{\FF_2}\ker T_J=L_J.
\]
另一方面，由定理 \ref{thm:conclusion-screen-entropy-audit-identity} 的证明，任意非空纤维 \(T_J^{-1}(y)\) 的基数都恰为
\[
2^{L_J},
\]
且在均匀先验下，条件分布 \(X\mid(Y_J=y)\) 在该纤维上均匀。

取任意可计算截面 \(s_J\)。则每条纤维都可唯一写为
\[
T_J^{-1}(y)=s_J(y)+\ker T_J
=
\left\{
s_J(y)+
\sum_{a_{J^c}}\varepsilon_{a_{J^c}}[\Sigma_{a_{J^c}}]:
\ \varepsilon_{a_{J^c}}\in\FF_2
\right\}.
\]
均匀分布在有限维 \(\FF_2\)-向量空间上，正对应于各坐标独立且均为公平 Bernoulli。故得到所示表示与超立方体识别。证毕。
\end{proof}

\begin{corollary}[后验相关结构的板层刚性]\label{cor:conclusion-coordinate-bundle-posterior-slab-rigidity}
在定理 \ref{thm:conclusion-coordinate-bundle-posterior-hypercube-factorization} 的设定下，任取两个顶维胞腔 \(C,C'\subset \mathcal Q_m\)。
\begin{enumerate}
  \item 若 \(C\) 与 \(C'\) 落在同一板层 \(\Sigma_{a_{J^c}}\) 中，则在条件 \(Y_J\) 下有
  \[
  X(C)-s_J(Y_J)(C)=X(C')-s_J(Y_J)(C').
  \]
  \item 若 \(C\subset \Sigma_{a_{J^c}}\)、\(C'\subset \Sigma_{b_{J^c}}\) 且 \(a_{J^c}\neq b_{J^c}\)，则这两个 centered bits 在条件 \(Y_J\) 下相互独立，且各自都是公平 Bernoulli。
\end{enumerate}
因此，内部坐标束屏幕的后验相关图具有严格二层结构：同层完全锁相，异层完全独立。
\end{corollary}
\begin{proof}
由定理 \ref{thm:conclusion-coordinate-bundle-posterior-hypercube-factorization}，
\[
X-s_J(Y_J)
=
\sum_{a_{J^c}}\varepsilon_{a_{J^c}}[\Sigma_{a_{J^c}}].
\]
若 \(C,C'\) 落在同一板层，则两者都只读取同一个系数 \(\varepsilon_{a_{J^c}}\)，故 centered bits 必相等。若二者落在不同板层，则对应不同的 \(\varepsilon\)-坐标，而这些坐标彼此独立且都服从公平 Bernoulli 分布。证毕。
\end{proof}

\begin{theorem}[外边界补全的 exact residual-rank 公式]\label{thm:conclusion-coordinate-bundle-outer-boundary-residual-rank}
\leanverified{paper\_conclusion\_coordinate\_bundle\_outer\_boundary\_residual\_rank}
以下假设 \(J\neq\varnothing\)。设 \(B\) 为任意外边界面族，并定义被击中的板层索引集
\[
\mathcal T_J(B):=
\left\{
a_{J^c}:\ 
B\cap \partial_{\mathrm{out}}\Sigma_{a_{J^c}}\neq\varnothing
\right\},
\]
其中 \(\partial_{\mathrm{out}}\Sigma_{a_{J^c}}\) 表示板层 \(\Sigma_{a_{J^c}}\) 与外边界相交的全部可见面。则有精确公式
\[
\rank_{\ZZ}\ker f_{S_J^{\mathrm{int}}\cup B}
=
L_J-|\mathcal T_J(B)|.
\]
\end{theorem}
\begin{proof}
由定理 \ref{thm:spg-internal-coordinate-bundle-screen-cost-closed-form}，
\(\Gamma_{S_J^{\mathrm{int}}}\) 的内部连通分量恰由各板层 \(\Sigma_{a_{J^c}}\) 给出，总数为 \(L_J\)；再加上外部顶点 \(\infty\)，初始连通分量数为
\[
c(\Gamma_{S_J^{\mathrm{int}}})=L_J+1.
\]

当加入一张外边界面 \(e\) 时，它只可能把 \(\infty\) 与某一条板层所在的内部连通分量连接起来；若该板层先前已被 \(B\) 中其他外边界面击中，则新加入的 \(e\) 不再进一步减少连通分量数。故加入整个 \(B\) 之后，图的连通分量数恰为
\[
c(\Gamma_{S_J^{\mathrm{int}}\cup B})=1+L_J-|\mathcal T_J(B)|.
\]
再由定理 \ref{thm:spg-screen-kernel-connected-components}，
\[
\rank_{\ZZ}\ker f_{S_J^{\mathrm{int}}\cup B}
=
c(\Gamma_{S_J^{\mathrm{int}}\cup B})-1
=
L_J-|\mathcal T_J(B)|.
\]
证毕。
\end{proof}

\begin{theorem}[外边界观测后的后验超立方体裁剪]\label{thm:conclusion-coordinate-bundle-posterior-hypercube-clipping}
以下仍假设 \(J\neq\varnothing\)。对任意外边界面族 \(B\)，记
\[
Y_{J,B}:=f_{S_J^{\mathrm{int}}\cup B}(X).
\]
则对每个 \(y\in \operatorname{im}Y_{J,B}\)，存在可计算截面
\[
s_{J,B}:\operatorname{im}Y_{J,B}\to C_n(\mathcal Q_m;\FF_2)
\]
使得
\[
X\mid (Y_{J,B}=y)
=
s_{J,B}(y)+
\sum_{a_{J^c}\notin \mathcal T_J(B)}\varepsilon_{a_{J^c}}[\Sigma_{a_{J^c}}],
\]
其中
\[
\bigl(\varepsilon_{a_{J^c}}\bigr)_{a_{J^c}\notin \mathcal T_J(B)}
\]
是 \(L_J-|\mathcal T_J(B)|\) 个相互独立的公平 Bernoulli 比特。特别地，
\[
X\mid (Y_{J,B}=y)\cong \FF_2^{\,L_J-|\mathcal T_J(B)|}
\]
作为均匀仿射超立方体。
\end{theorem}
\begin{proof}
由定理 \ref{thm:conclusion-coordinate-bundle-posterior-hypercube-factorization}，对每个
\[
y_0\in \operatorname{im}T_J
\]
都有
\[
X\mid (Y_J=y_0)
=
s_J(y_0)+\sum_{a_{J^c}}\varepsilon_{a_{J^c}}[\Sigma_{a_{J^c}}],
\]
其中全部 \(\varepsilon_{a_{J^c}}\) 相互独立且公平。

另一方面，由定理 \ref{thm:conclusion-coordinate-bundle-outer-boundary-residual-rank} 与推论 \ref{cor:conclusion-coordinate-bundle-single-face-strict-zero-one-law}，加入 \(B\) 之后，恰有被命中板层
\[
a_{J^c}\in \mathcal T_J(B)
\]
对应的自由坐标被逐一固定，而未命中板层的自由坐标保持独立且不受耦合。故每条后验纤维都可唯一写成
\[
s_{J,B}(y)+
\sum_{a_{J^c}\notin \mathcal T_J(B)}\varepsilon_{a_{J^c}}[\Sigma_{a_{J^c}}],
\]
其中 \(s_{J,B}\) 可通过逐层读出被固定坐标后显式构造。由于剩余坐标仍在 \(\FF_2\) 上均匀独立，便得到所示仿射超立方体识别。证毕。
\end{proof}

\begin{corollary}[坐标束屏幕的最小外边界精确化阈值]\label{cor:conclusion-coordinate-bundle-minimal-outer-boundary-threshold}
以下仍假设 \(J\neq\varnothing\)。使
\[
f_{S_J^{\mathrm{int}}\cup B}
\]
成为单射所需的最小外边界面数恰为
\[
L_J=2^{m(n-s)}.
\]
等价地，必须且只需从每一条板层 \(\Sigma_{a_{J^c}}\) 的外边界中至少选取一个面。
\end{corollary}
\begin{proof}
由定理 \ref{thm:conclusion-coordinate-bundle-outer-boundary-residual-rank}，
\[
f_{S_J^{\mathrm{int}}\cup B}\ \text{单射}
\iff
\rank_{\ZZ}\ker f_{S_J^{\mathrm{int}}\cup B}=0
\iff
|\mathcal T_J(B)|=L_J.
\]
另一方面，由 \(J\neq\varnothing\) 与推论 \ref{cor:spg-coordinate-bundle-minimal-boundary-closure} 的论证，每条板层都含有非空外边界，而每张外边界面至多只会击中一条板层。因此至少需要 \(L_J\) 张面；反过来，若每层各取一面，则 \(|\mathcal T_J(B)|=L_J\)，故核秩为零。证毕。
\end{proof}

\begin{theorem}[带权外边界精确化的分层最优化公式]\label{thm:conclusion-coordinate-bundle-weighted-boundary-exactification}
以下仍假设 \(J\neq\varnothing\)。赋予每个外边界面 \(e\) 一个非负权重 \(\omega(e)\)。定义 exactifying family 的最小总权重
\[
W_J^{\min}:=
\min\left\{
\sum_{e\in B}\omega(e):
\ f_{S_J^{\mathrm{int}}\cup B}\ \text{单射}
\right\}.
\]
则有显式公式
\[
W_J^{\min}
=
\sum_{a_{J^c}\in\{0,\dots,2^m-1\}^{J^c}}
\min_{e\in \partial_{\mathrm{out}}\Sigma_{a_{J^c}}}\omega(e).
\]
因此，内部坐标束屏幕的带权精确化并非一般图优化问题，而是逐板层完全可分离的线性选择问题。
\end{theorem}
\begin{proof}
由推论 \ref{cor:conclusion-coordinate-bundle-minimal-outer-boundary-threshold}，精确化当且仅当每条板层至少命中一次外边界。又由于一张外边界面只隶属于某一条固定板层的外边界，故不同板层之间的约束完全解耦。于是全局最优值只能等于：对每条板层分别取其外边界上权重最小的一张面，然后把这些最小值求和。证毕。
\end{proof}

\begin{corollary}[单面边际收益的 strict \texorpdfstring{$0$--$1$}{0-1} 律]\label{cor:conclusion-coordinate-bundle-single-face-strict-zero-one-law}
\leanverified{paper\_conclusion\_coordinate\_bundle\_single\_face\_strict\_zero\_one\_law}
以下仍假设 \(J\neq\varnothing\)。对任意外边界面 \(e\)，设它属于板层 \(\Sigma_{a_{J^c}(e)}\)。则对任意已有补全面集 \(B\)，都有
\[
\rank_{\ZZ}\ker f_{S_J^{\mathrm{int}}\cup B}
-
\rank_{\ZZ}\ker f_{S_J^{\mathrm{int}}\cup(B\cup\{e\})}
=
\mathbf 1_{\{a_{J^c}(e)\notin \mathcal T_J(B)\}}.
\]
若再记
\[
Y_{J,B}:=f_{S_J^{\mathrm{int}}\cup B}(X),
\]
则条件熵下降也满足同一 \(0\text{--}1\) 律：
\[
H(X\mid Y_{J,B})-H(X\mid Y_{J,B\cup\{e\}})
=
\mathbf 1_{\{a_{J^c}(e)\notin \mathcal T_J(B)\}}.
\]
\end{corollary}
\begin{proof}
由定理 \ref{thm:conclusion-coordinate-bundle-outer-boundary-residual-rank}，
\[
\rank_{\ZZ}\ker f_{S_J^{\mathrm{int}}\cup B}=L_J-|\mathcal T_J(B)|.
\]
若 \(a_{J^c}(e)\in\mathcal T_J(B)\)，则加入 \(e\) 不改变 \(\mathcal T_J(B)\)，差值为 \(0\)。若 \(a_{J^c}(e)\notin\mathcal T_J(B)\)，则
\[
|\mathcal T_J(B\cup\{e\})|=|\mathcal T_J(B)|+1,
\]
差值为 \(1\)。这就证明了第一式。

第二式由定理 \ref{thm:conclusion-screen-entropy-audit-identity} 立即得到，因为对任意屏幕 \(S\)，都有
\[
H(X\mid f_S(X))=\rank_{\ZZ}\ker f_S.
\]
把 \(S=S_J^{\mathrm{int}}\cup B\) 与 \(S=S_J^{\mathrm{int}}\cup(B\cup\{e\})\) 代入即可。证毕。
\end{proof}

\begin{theorem}[最优线性补全的 torsor 分类]\label{thm:conclusion-coordinate-bundle-optimal-linear-completion-torsor}
固定一个线性核坐标映射
\[
\kappa_J:C_n(\mathcal Q_m;\FF_2)\to \FF_2^{L_J}
\]
与一个线性截面
\[
s_J:\operatorname{im}T_J\to C_n(\mathcal Q_m;\FF_2)
\]
满足：
\[
\kappa_J|_{\ker T_J}:\ker T_J\xrightarrow{\sim}\FF_2^{L_J}
\quad\text{为同构，}\qquad
\kappa_J\circ s_J=0.
\]
设
\[
L:C_n(\mathcal Q_m;\FF_2)\to \FF_2^{L_J}
\]
是任一线性映射，并且联合映射
\[
x\longmapsto \bigl(T_Jx,Lx\bigr)
\]
为单射。则存在唯一
\[
U\in \mathrm{GL}_{L_J}(\FF_2),
\qquad
R\in \operatorname{Hom}\!\bigl(\operatorname{im}T_J,\FF_2^{L_J}\bigr)
\]
使得
\[
L=U\kappa_J+R\,T_J.
\]
换言之，全部最优线性补全恰组成一个
\[
\operatorname{Hom}\!\bigl(\operatorname{im}T_J,\FF_2^{L_J}\bigr)\rtimes \mathrm{GL}_{L_J}(\FF_2)
\]
上的 torsor。
\end{theorem}
\begin{proof}
由定理 \ref{thm:conclusion-relative-linear-holography-exact-supplement-dimension}，\(L_J=\dim_{\FF_2}\ker T_J\) 正是最小线性补全维数。若联合映射 \((T_J,L)\) 单射，则
\[
L|_{\ker T_J}:\ker T_J\to \FF_2^{L_J}
\]
必为同维线性同构。另一方面，
\[
\kappa_J|_{\ker T_J}:\ker T_J\to \FF_2^{L_J}
\]
也是同构，因此存在唯一
\[
U\in\mathrm{GL}_{L_J}(\FF_2)
\]
使
\[
L|_{\ker T_J}=U\,\kappa_J|_{\ker T_J}.
\]

再对任意 \(y\in \operatorname{im}T_J\) 定义
\[
R(y):=L(s_J(y)).
\]
由于 \(s_J\) 与 \(L\) 都线性，故 \(R\) 线性。任取 \(x\in C_n(\mathcal Q_m;\FF_2)\)，令
\[
k:=x-s_J(T_Jx)\in\ker T_J.
\]
则
\[
\begin{aligned}
L(x)
&=L\bigl(s_J(T_Jx)\bigr)+L(k)\\
&=R(T_Jx)+U\,\kappa_J(k)\\
&=R(T_Jx)+U\,\kappa_J(x),
\end{aligned}
\]
其中最后一步用到了 \(\kappa_J\circ s_J=0\)。故
\[
L=U\kappa_J+R\,T_J.
\]
唯一性由 \(U\) 在 \(\ker T_J\) 上的唯一确定性以及 \(R(y)=L(s_J(y))\) 的定义立即得到。证毕。
\end{proof}

\begin{theorem}[Shannon、audit、linear holography 与 Kolmogorov 的四重复杂度一致坍缩]\label{thm:conclusion-coordinate-bundle-fourfold-complexity-collapse}
对内部坐标束屏幕 \(S_J^{\mathrm{int}}\)，以下四个量在本模型中完全一致，其中 Kolmogorov 项只差 \(O(1)\)：
\[
H(X\mid Y_J),
\]
\[
\min\Bigl\{
b:\exists\ \text{\(b\)-bit exact audit completion for }T_J
\Bigr\},
\]
\[
\min\Bigl\{
L:\exists\ L\ \text{个标量线性泛函使 }x\mapsto (T_Jx,\ell_1(x),\dots,\ell_L(x))\ \text{单射}
\Bigr\},
\]
\[
\max_{y\in\operatorname{im}T_J}\ \max_{x:T_Jx=y} K(x\mid y,m,n,J).
\]
更精确地，前三者都恰等于
\[
L_J=2^{m(n-s)},
\]
而第四者满足
\[
\max_{y\in\operatorname{im}T_J}\ \max_{x:T_Jx=y} K(x\mid y,m,n,J)
=
L_J+O(1).
\]
\end{theorem}
\begin{proof}
由定理 \ref{thm:conclusion-screen-entropy-audit-identity}，
\[
H(X\mid Y_J)=\rank_{\ZZ}\ker T_J.
\]
再由定理 \ref{thm:conclusion-coordinate-bundle-kernel-slab-decomposition}，
\[
\rank_{\ZZ}\ker T_J=L_J.
\]
故第一项等于 \(L_J\)。

同样由定理 \ref{thm:conclusion-screen-entropy-audit-identity}，最小二进审计补全位数也恰等于 \(\rank_{\ZZ}\ker T_J=L_J\)。这给出第二项。

第三项则由定理 \ref{thm:conclusion-relative-linear-holography-exact-supplement-dimension} 直接给出：使 \(T_J\) 变成单射所需的最少标量线性泛函个数正是
\[
2^{m(n-s)}=L_J.
\]

最后，对 Kolmogorov 项，推论 \ref{cor:conclusion-screen-max-conditional-complexity-equals-rank} 作用于屏幕 \(S_J^{\mathrm{int}}\) 给出：对任意
\[
y\in\operatorname{im}T_J,
\]
都有
\[
\max_{x:T_Jx=y}K(x\mid y,m,n,J)=\rank_{\ZZ}\ker T_J+O(1)=L_J+O(1).
\]
对 \(y\) 再取最大值仍得同样量级。证毕。
\end{proof}

\begin{theorem}[坐标方向的 exact \texorpdfstring{$m$}{m}-bit drop law]\label{thm:conclusion-coordinate-bundle-exact-mbit-drop-law}
\leanverified{paper\_conclusion\_coordinate\_bundle\_exact\_mbit\_drop\_law}
定义隐藏体积函数
\[
\Delta_m(J):=\rank_{\ZZ}\ker f_{S_J^{\mathrm{int}}}=2^{m(n-|J|)}.
\]
则对任意 \(j\notin J\)，有精确比率律
\[
\frac{\Delta_m(J\cup\{j\})}{\Delta_m(J)}=2^{-m},
\]
等价地，
\[
\log_2\Delta_m(J)-\log_2\Delta_m(J\cup\{j\})=m.
\]
相应地，可见秩增量满足
\[
\rank_{\QQ}\operatorname{im}T_{J\cup\{j\}}
-
\rank_{\QQ}\operatorname{im}T_J
=
(2^m-1)\,2^{m(n-|J|-1)}.
\]
因此，每增加一个坐标方向，可见系统都会无条件地消灭恰好 \(m\) 比特的隐藏量；这一降幅与所加入的方向 \(j\) 本身无关。
\end{theorem}
\begin{proof}
由定理 \ref{thm:conclusion-coordinate-bundle-logmodularity-visible-rank}，
\[
\Delta_m(J)=2^{m(n-|J|)},
\qquad
\rank_{\QQ}\operatorname{im}T_J=2^{mn}-2^{m(n-|J|)}.
\]
故
\[
\frac{\Delta_m(J\cup\{j\})}{\Delta_m(J)}
=
2^{m(n-|J|-1)-m(n-|J|)}
=
2^{-m},
\]
取 \(\log_2\) 即得第二式。最后一式直接由可见秩公式相减得到：
\[
\bigl(2^{mn}-2^{m(n-|J|-1)}\bigr)-\bigl(2^{mn}-2^{m(n-|J|)}\bigr)
=(2^m-1)\,2^{m(n-|J|-1)}.
\]
证毕。
\end{proof}

\begin{theorem}[坐标信息游戏的 Möbius 原子性与 Shapley 均分律]\label{thm:conclusion-coordinate-bundle-mobius-atomicity-shapley-law}
\leanverified{paper\_conclusion\_coordinate\_bundle\_mobius\_atomicity\_shapley\_law}
定义
\[
h_m(J):=\log_2\Delta_m(J)=m(n-|J|).
\]
把 \(h_m\) 视为布尔格 \(2^{[n]}\) 上的集合函数，并记其 Möbius 反演为
\[
\widehat h_m(T):=
\sum_{U\subseteq T}(-1)^{|T|-|U|}h_m(U).
\]
则
\[
\widehat h_m(\varnothing)=mn,
\qquad
\widehat h_m(\{j\})=-m\ \ (1\le j\le n),
\qquad
\widehat h_m(T)=0\ \ (|T|\ge 2).
\]
等价地，全部二阶及以上交互项都严格消失。

于是对收益函数
\[
v_m(J):=h_m(\varnothing)-h_m(J)=m|J|,
\]
每个坐标方向的 Shapley 值都恰为
\[
\operatorname{Sh}_j(v_m)=m.
\]
\end{theorem}
\begin{proof}
由定理 \ref{thm:conclusion-coordinate-bundle-exact-mbit-drop-law}，
\[
h_m(J)=m(n-|J|)=mn-m|J|.
\]
因此 \(h_m\) 是布尔格指标向量 \((\mathbf 1_{j\in J})_{j=1}^{n}\) 的仿射函数。对这样的函数，Möbius 反演只可能在 \(\varnothing\) 与单点集上非零。直接计算得
\[
\widehat h_m(\varnothing)=h_m(\varnothing)=mn,
\]
以及
\[
\widehat h_m(\{j\})=h_m(\{j\})-h_m(\varnothing)=-m.
\]
若 \(|T|\ge 2\)，则在
\[
\sum_{U\subseteq T}(-1)^{|T|-|U|}\bigl(mn-m|U|\bigr)
\]
中，常数项与一次项都发生完全抵消，故
\[
\widehat h_m(T)=0.
\]

对收益函数 \(v_m(J)=m|J|\)，任一坐标 \(j\) 的边际贡献
\[
v_m(J\cup\{j\})-v_m(J)=m
\]
与加入顺序及背景集合 \(J\) 无关。故 Shapley 平均值不需再作非平凡加权，直接等于该常数 \(m\)。证毕。
\end{proof}

\noindent
由此，内部坐标束屏幕在结论层形成了一条完全显式的闭链：板层示性链不只给出核的整数基，而且把每条后验纤维精确分裂为一个 \(L_J\) 维独立超立方体；一旦加入外边界观测，这个后验超立方体又仅按被命中板层逐坐标裁剪，而不会生成任何跨层耦合；外边界精确化则不再是一般的增秩过程，而是“每层命中一次即可”的残余秩削减问题，其带权版本逐板层完全分离，单面边际收益严格服从 \(0\text{--}1\) 律；再往上线性化，全部最优补全恰组成一个显式的 affine torsor；最终，Shannon 熵、最小审计位数、最小线性补全维数与最大条件 Kolmogorov 容量在同一整数 \(L_J\) 上一致坍缩，而坐标方向的协同结构又在布尔格 Möbius 反演下退化为纯单体原子，并以 Shapley 值 \(m\) 的形式均匀分配。于是，内部坐标束全息并不隐藏任何更高阶交互：其全部真实复杂度都精确锁定在板层个数与单坐标原子之上。

\endinput
